\section{Conclusion and Future Work}
\label{sec:conclusion}

We have presented a semantic account of data provenance as a form of automatic differentiation, representing
the dependency of outputs on inputs as a derivative whose Jacobian is a matrix over a commutative semiring.
Traditional automatic differentiation, and various program analysis techniques such as \GPS, arise as
instances. The model interprets an expressive higher-order language for querying and manipulating data, and
reveals new applications such as approximation by intervals~(\secref{approx-as-tangents:perturbation-bounds}).
It also opens up an opportunity to integrate traditional symbolic data provenance techniques with the AD-based
methods used in explainable AI, such as the saliency maps of Grad-CAM~\cite{selvaraju20}. Our categorical
approach admits a modular construction of the model, and the use of general theorems, such as Fiore and
Simpson's \thmref{glr-definability}, to prove properties of the interpretation. We have focused on
constructions that enable an executable implementation in Agda.

\paragraph{Alternative semiring models}

The choice of scalar semiring is the principal degree of freedom in our framework, and choices
beyond those studied here seem worth pursuing. The perturbation-bound semirings of
\secref{examples:perturbation-bounds} point towards metric approximation: a Lawvere metric space is a
category enriched over the min-plus quantale, whose operations are those of our absolute-perturbation
semiring, so metric spaces provide a native notion of approximation already partly within reach of the
framework. Relating this to \citet{edalat-heckmann98}'s computational model of metric spaces is one promising
direction. A second is to replace the setoids carrying values with sets equipped with a semiring-valued
equality relation, over a partially ordered semiring with a residual for its multiplication.

% \paragraph{Categorical Models of Differentation} We have already identified \conref{tangent-stable-fns} and
% \conref{tangent-fam} as potential connections to categorical models of differentation.

\paragraph{Connection to sequentiality and stable functions}

In our model the forward derivative carried by a morphism of $\Fam(\SemiMod(S))$ is supplied as part
of the interpretation, not determined by the underlying function, so it need not reflect how that function
depends on its input. Berry's
\emph{stable functions}~\cite{berry79,berry82}, by contrast, come equipped with an intrinsic forward and
backward derivative. Stability was proposed as a refinement of domain theory aimed at capturing the
intensional behaviour of sequential programs, imposing constraints on how functions preserve bounded meets of
approximants; although stable functions do not fully characterise sequentiality, because they admit
$\mathrm{gustave}$-style counterexamples, they remain an appropriate notion for studying the sensitivity of a
program to partial data at a specific point~\cite{amadio-curien}. Paul Taylor's characterisation of stable
functions via local Galois connections on principal downsets provides a semantic underpinning for the reverse
maps of Galois slicing~\cite{taylor99}. The stable-function story has a tight connection to Galois slicing but a looser one to automatic
differentiation, so we intend to develop it in separate work, relating Galois slicing to stable functions
rather than to differentiation. Such a development would also clarify the links to Differential Linear
Logic~\cite{ehrhard_differential_2006} and to Crubill\'e's identification of stable functions with power
series~\cite{crubille18}, relating the qualitative derivatives studied here to the analytic derivatives of
that setting.

\paragraph{Shape Approximation}

Lists are the only recursive datatype in our source language, so supporting general inductive and
coinductive types is important future work, perhaps adapting \posscite{nunes2023} automatic differentiation
for $\mu\nu$-polynomial functors. Richer datatypes bring the question of \emph{shape approximation}:
approximating not only the values at the positions of a structure but the structure itself, such as the
length of a list or the depth of a tree. Shape approximation features in earlier work on Galois
slicing~\cite{perera12a,perera16d,ricciotti17}, whose approximation lattices represent partially erased data,
whereas our model holds the shape of each input fixed. Our Jacobians are spans, and polynomial functors are
the shape-dependent generalisation of spans~\cite{gambino-kock}; the finite polynomials form the Lawvere
theory for commutative semirings, the algebra of our scalars. The corresponding datatypes are containers,
whose derivative is the type of one-hole contexts, carrying a notion of linear map and a chain rule for both
inductive and coinductive types~\cite{abbott05}. This already gives a set-theoretic account of shape
approximation; weighting positions with our semiring scalars would extend our analyses to varying shapes.

\paragraph{Recursion and Partiality}  This work treats only total programs, and makes no use of directed completeness or similar
properties of domains. Extending to recursion and partiality would introduce a notion of undefinedness, and
hence a definedness order to be reconciled with any approximation order on the same data. Non-termination and
infinite data are both partiality phenomena, approximating unbounded computation by finite observation, so
potentially the finite prefixes of an infinite or non-terminating computation would fall under the same
extension.

\paragraph{Source-To-Source Translation Techniques}

% To formalise this, we'd need to... See also $\Fam(\PSh_{\CMon}(\namedcat{Syn}))$?

An interesting alternative to the denotational approach presented here, and to the trace-based approaches used
in some implementations of provenance, would be to develop a
source-to-source transformation, in direct analogy with
the CHAD approach to automatic differentiation \cite{vakar22,nunes2023}. In their approach, forward and
reverse-mode AD are implemented as compositional transformations on source code, guided by a universal
property: they arise as the unique structure-preserving functors from the source language to a suitably
structured target language formalised as a Grothendieck construction. Adapting this to our setting
would allow the analysis to be ``compiled in'', avoiding the need for
a custom interpreter and potentially exposing opportunities for optimisation.
